\section{Place four independent quarter-rounds in vector lanes}
A word-format MSA register contains four 32-bit lanes. Place the first state row
in \texttt{w0}, the second in \texttt{w1}, the third in \texttt{w2}, and the fourth
in \texttt{w3}. Number the lanes zero through three in increasing memory order.
A word-vector load at state offset sixteen therefore gives \texttt{w1} the key
words at indices four through seven.

In this arrangement lane zero across the four registers holds state indices
\((0,4,8,12)\); lane one holds \((1,5,9,13)\). These are different quarter-rounds
with no shared words. The same sequence of vector operations can update all four
lanes independently. Within each lane, the dependent order from the scalar
quarter-round still applies.

Preserve original rows in \texttt{w4} through \texttt{w7} before changing the working
rows. The vector version has enough registers to retain all sixteen original words
without writing a second caller buffer. The instruction \texttt{move.v} copies a
complete register; it does not change the row's lane order.
\Needspace{9\baselineskip}
\MARSCodeLabel{Load and preserve the first two rows.}
\begin{minted}[samepage=true]{mars-idea}
ld.w $w0, 0($a0)
move.v $w4, $w0
ld.w $w1, 16($a0)
move.v $w5, $w1
\end{minted}

Use \texttt{addv.w} for independent modulo-\(2^{32}\) lane sums. A carry out of lane
zero does not enter lane one. \texttt{xor.v} and \texttt{or.v} are bitwise operations
over the vector; because they do not move bits between positions, they also respect
these word boundaries. A lane rotation needs word shifts, not a shift of the full
128-bit register.
\Needspace{10\baselineskip}
\MARSCodeLabel{Perform the first quarter-round group in all four lanes.}
\begin{minted}[samepage=true]{mars-idea}
addv.w $w0, $w0, $w1
xor.v $w3, $w3, $w0
slli.w $w8, $w3, 16
srli.w $w9, $w3, 16
or.v $w3, $w8, $w9
\end{minted}

The two temporary registers hold the two parts of each lane's rotation. Continue
with the three remaining groups using left amounts twelve, eight and seven and
complementary logical right amounts twenty, twenty-four and twenty-five.
Completing these groups finishes one column round for the entire state. It does not
finish a complete block.
